% =============================================================================
% БИЛЕТ 14
% =============================================================================

\ticket{14}{}

\question{Достаточное условие экстремума ФНП}

\begin{defn}[Стационарная точка]
  Точка $x^0$ называется \emph{стационарной}, если все частные производные $f$ в ней равны нулю:
  $\dfrac{\partial f}{\partial x_1}(x^0) = \ldots = \dfrac{\partial f}{\partial x_n}(x^0) = 0$.
\end{defn}

\begin{defn}[Матрица Гессе и второй дифференциал]
  Пусть $f \in C^2(\mathbb{R}^n, B_\delta(x^0))$. Второй дифференциал:
  \[
    d^2_{x^0} f(dx) = \sum_{i=1}^n\sum_{j=1}^n \frac{\partial^2 f}{\partial x_i\,\partial x_j}(x^0)\,dx_i\,dx_j.
  \]
  Матрица $H_f(x^0) = \left(\dfrac{\partial^2 f}{\partial x_i\,\partial x_j}(x^0)\right)_{i,j=1}^n$
  называется \emph{матрицей Гессе}. Квадратичная форма
  $K(h) = \sum_{i,j} a_{ij} h_i h_j$, $h_i = dx_i$, $a_{ij} = a_{ji}$.
\end{defn}

\begin{defn}[Знакоопределённость формы]
  Форма $K$ \emph{положительно определена}, если $\forall h \neq 0: K(h) > 0$;
  \emph{отрицательно определена}, если $K(h) < 0$ при $h \neq 0$;
  \emph{знаконеопределена}, если принимает и положительные, и отрицательные значения.
\end{defn}

\begin{theorem}[Критерий Сильвестра]
  $K$ положительно определена $\Leftrightarrow$ все угловые миноры матрицы $H_f$ положительны.
  $K$ отрицательно определена $\Leftrightarrow$ угловые миноры чередуются: $\Delta_1 < 0$, $\Delta_2 > 0$, \ldots
\end{theorem}

\begin{lemma}
  Если $K$ положительно определена, то $\exists m > 0: K(h) \geq m\,\|h\|^2$.
  \begin{proof}
    $K$ непрерывна; ограничение на единичную сферу $\{h: \|h\|=1\}$ (компакт) достигает минимума $m > 0$.
    Для произвольного $h \neq 0$ подставляем $\tilde{h} = h/\|h\|$:
    $K(h) = \|h\|^2 K(\tilde{h}) \geq m\,\|h\|^2$. \qedhere
  \end{proof}
\end{lemma}

\begin{theorem}[Достаточное условие экстремума]
  Пусть $x^0$ --- стационарная точка $f \in C^2$.
  \begin{enumerate}
    \item Если $d^2_{x^0}f(dx)$ положительно определена, то $x^0$ --- строгий локальный минимум.
    \item Если $d^2_{x^0}f(dx)$ отрицательно определена, то $x^0$ --- строгий локальный максимум.
    \item Если $d^2_{x^0}f(dx)$ знаконеопределена, то экстремума нет.
  \end{enumerate}
  \begin{proof}
    По формуле Тейлора:
    \[
      f(x) - f(x^0) = \underbrace{d_{x^0}f}_{=0} + \frac{d^2_{x^0}f(dx)}{2} + O(\|x-x^0\|^2).
    \]
    \textbf{Пункты 1--2.} По лемме $d^2_{x^0}f(dx) \geq m\,\|dx\|^2$.
    Остаток $O(\|x-x^0\|^2)$ можно сделать $< \tfrac{m}{2}\|x-x^0\|^2$ в достаточно малом шаре.
    Тогда $f(x) - f(x^0) \geq \tfrac{m}{4}\|x-x^0\|^2 > 0$ --- строгий минимум.
    Для пункта 2 аналогично.

    \textbf{Пункт 3.} Существуют направления $h_+$, $h_-$ такие, что $d^2_{x^0}f(h_\pm) \gtrless 0$.
    Вдоль каждого из них знак $f(x) - f(x^0)$ не меняется (малое $o$ не меняет знак)
    и различается для $h_+$ и $h_-$ --- экстремума нет. \qedhere
  \end{proof}
\end{theorem}

\question{Интегрируемость монотонной, непрерывной и кусочно-непрерывной функции как следствие критерия Лебега}

\begin{defn}[Мера нуль по Лебегу]
  Множество $E \subset \mathbb{R}$ имеет \emph{лебегову меру нуль}, если
  $\forall\varepsilon>0\;\exists$ конечное или счётное покрытие интервалами $\{(a_i, b_i)\}$:
  $E \subset \bigcup (a_i, b_i)$ и $\sum |b_i - a_i| < \varepsilon$.
\end{defn}

\begin{lemma}
  Любое счётное множество имеет лебегову меру нуль.
  \begin{proof}
    $E = \{x_k\}_{k=1}^\infty$. Накрываем $x_k$ интервалом длины $\varepsilon/2^{k+1}$;
    суммарная длина $\leq \varepsilon$. \qedhere
  \end{proof}
\end{lemma}

\begin{lemma}
  Счётное объединение множеств меры нуль имеет меру нуль.
  \begin{proof}
    Пусть $E_i$ имеет меру нуль. Накрываем $E_i$ с суммарной длиной $< \varepsilon/2^i$;
    объединив покрытия, получим суммарную длину $< \varepsilon$. \qedhere
  \end{proof}
\end{lemma}

\begin{theorem}[Критерий Лебега]
  $f: [a,b] \to \mathbb{R}$ интегрируема по Риману тогда и только тогда, когда:
  \begin{enumerate}
    \item $f$ ограничена на $[a,b]$;
    \item множество $D[f]$ точек разрыва $f$ имеет лебегову меру нуль.
  \end{enumerate}
\end{theorem}

\begin{corollary}[Непрерывная $\Rightarrow$ интегрируема]
  Если $f \in C([a,b])$, то $f \in R([a,b])$.
  \begin{proof}
    $f$ непрерывна на $[a,b]$ $\Rightarrow$ ограничена (теорема Вейерштрасса).
    Точек разрыва нет $\Rightarrow$ $D[f] = \emptyset$ имеет меру нуль. \qedhere
  \end{proof}
\end{corollary}

\begin{corollary}[Монотонная $\Rightarrow$ интегрируема]
  Если $f$ монотонна на $[a,b]$, то $f \in R([a,b])$.
  \begin{proof}
    Монотонная $f$ ограничена: $f(a) \leq f(x) \leq f(b)$ (или наоборот).
    Разрывы монотонной функции 1-го рода; их счётно, поэтому $D[f]$ счётно
    $\Rightarrow$ мера нуль. \qedhere
  \end{proof}
\end{corollary}

\begin{corollary}[Кусочно-непрерывная $\Rightarrow$ интегрируема]
  Если $f$ кусочно-непрерывна на $[a,b]$, то $f \in R([a,b])$.
  \begin{proof}
    $f$ ограничена. Точки разрыва --- конечный набор точек из разбиения (их конечно)
    $\Rightarrow$ $D[f]$ конечно $\Rightarrow$ мера нуль. \qedhere
  \end{proof}
\end{corollary}
